\uuid{gWfm}
\exo7id{4363}
\titre{exo7 4363}
\auteur{quercia}
\organisation{exo7}
\datecreate{2010-03-12}
\isIndication{false}
\isCorrection{true}
\chapitre{Intégration}
\sousChapitre{Intégrale de Riemann dépendant d'un paramètre}
\cpge{mpsi-mp2i.nombres-complexes.c, pcsi.nombres-complexes.c, ptsi.nombres-complexes.c, mpsi-mp2i.nombres-reels-suites.c, pcsi.nombres-reels-suites.b, ptsi.nombres-reels-suites.b, mpsi-mp2i.nombres-reels-suites.d, pcsi.nombres-reels-suites.c, ptsi.nombres-reels-suites.c, mpsi-mp2i.procedes-sommatoires-discrets.a, pcsi.series-numeriques.a, ptsi.series-numeriques.a, bcpst1.nombres-complexes.a, bcpst1.suites-reelles-usuelles, bcpst1.suites-reelles, bcpst1.integration-segment.a, mp-mpi.topologie-espaces-normes.b, pc.espaces-vectoriels-normes.b, psi.espaces-vectoriels-normes.b, mp-mpi.series-numeriques-vectorielles.a, pt.series-numeriques.b, mp-mpi.suites-series-fonctions.c, bcpst2.series-reelles}
\module{Analyse}
\niveau{L2}
\difficulte{}

\contenu{
\question{
Soit $\alpha \in \R$. Trouver la limite de $u_n=\sum_{k=1}^n \frac{\sin (k\alpha)}{n+k}\cdotp$
}
\reponse{
Si $\alpha\in2\pi\,\Z$ alors $u_n=0$ pour tout~$n$.

Sinon, $u_n = \Im\Bigl( \int_{t=0}^1 \sum_{k=1}^nt^{n+k-1}e^{ik\alpha}\,d t\Bigr)
=\Im\Bigl( \int_{t=0}^1\frac{t^ne^{i\alpha}-t^{2n}e^{i(n+1)\alpha}}{1-te^{i\alpha}}\,d t\Bigr)
\to 0$ (lorsque $n\to\infty$) par convergence dominée.
}
}
